\section{The Reader As Final Station}

The reader could not have been the first station. At the beginning, a reader
who simply chose an orientation would have been an unlicensed authority. The
grammar had to earn difference, identity, finite forcing, the positive
invariant, baseline-relative sign, boundary trace, outside orientation,
quotient selection, holonomy, native count, and gauge partition. Only after
those obligations are in place can the reader stand at the end without
becoming a new primitive.

The final station is therefore not a person outside the grammar. It is the
role in which the grammar's predicates are read. A predicate is a question
with a permitted form: did this distinguish, did this fall inside the
allowed magnitude band, did this rise to the next admissible reading, did
this pool, did this select, did this close. The reader is addressed by those
questions because the grammar has no further interior place to hide them.

This address is constrained. The reader may orient a convention once the
boundary has made orientation lawful. The reader may select a quotient once
the observation has specified its equivalence relation. The reader may read a
holonomy sign once the loop has closed. The reader may count native matter
and admitted antimatter once the baseline and gauge partition have been
fixed. The reader may stop or continue explanation once a predicate changes
or fails to change the boundary signal. The reader may not invent a new
invariant, erase an old residue, or turn an unchanged presentation into a
new measurement.

This makes the final station both modest and decisive. It is modest because
it does not own the world. It receives questions formed by the grammar. It
is decisive because, once those questions have been formed, no further
interior process can answer them in the reader's place. The station reads
whether the final distinction is present. If the distinction is present, the
grammar permits the next step. If not, the grammar forbids another layer as
idle repetition.

The first question is distinguishability. Did this reading make a difference
that the selected grammar can recover? If yes, the reader may carry the
difference forward. If no, the presentations are identified for this
reading. This is the first permission from the foundation returned at the
end of the chain. The book began by allowing a difference to be read. The
final station asks whether a proposed last difference is actually readable.

The second question is magnitude. A reading may be too small to cross the
noise floor, inside the admissible band, or too large for the current gauge
to carry without changing the problem. The reader does not need an infinite
scale of all possible sizes. The reader needs the band the grammar has made
available. Below the band is not yet a carried distinction. Inside the band
is admissible. Above the band forces a new reading or an obstruction. The
station reads which case holds.

The third question is rise. A strict order is not given at the foundation.
It is found when one reading is below another in the permitted preorder and
the reverse reading is not permitted. The reader may therefore ask whether a
proposed step actually rises to a new station or merely renames the old one.
If it rises, the grammar carries a new level. If it does not, the attempted
order collapses into equivalence or neutrality.

The gauge questions then become reader-facing. Did this gate pool? Did it
select? Did the gauge close? The reader does not decide these by preference.
The gate has already specified what pooling and selecting mean for the
reading. The reader answers by applying the predicate. If every gate is
accounted for, the gauge closes. If a gate escapes both sides, the claimed
partition has failed. The final station is the audit point at which the
finite gauge is actually read as finite.

This is also where the matter/native asymmetry becomes an address rather
than a slogan. The reader is asked: under the native baseline, does any gate
select the positive sign? The answer is no. Under the tilted comparison, is
there a closed selection that reads the positive sign? The answer may be yes
when the gate has cranked. The reader is not choosing matter over antimatter.
The reader is answering the predicates that the baseline and gauge have
forced.

The final station also binds relativity to responsibility. Because every
relative reading depends on a selected quotient, the reader must say which
differences have been identified and which residue is carried. Because every
sign depends on orientation, the reader must say which baseline or closed
loop has oriented the sign. Because every finite count depends on a gauge,
the reader must say how the gates partition. Relativity is not escape from
accountability. It is accountability to the selected grammar of the reading.

This is why the chapter can call the reader the final apparatus only in the
grammatical sense. The final station is not a machine that produces the
chapter's result. It is the last place at which the grammar's questions can be answered.
The predicates turn back on the reader because the reader is the one
authorized to select, orient, read, stop, and continue after the boundary has
made those acts lawful.

The chapter therefore closes the forward path. Relative velocity was read as
quotient selection. The positron was read as closed-loop holonomy sign.
Native asymmetry was counted without importing antimatter. The thirty-six
gates partitioned into pooling and selecting moves. The predicates exposed
by those moves turned back on the reader. What remains is not another
forward addition. The next chapter can begin the backward walk, returning
through the tower to see how the loop closes around the single needle of
identity.
